\documentclass[10pt,letterpaper,sans]{moderncv}
\moderncvstyle{banking}
\moderncvcolor{blue}

\usepackage[scale=0.84]{geometry}
\nopagenumbers{}
\setlength{\hintscolumnwidth}{2.35cm}

\name{Krishnaa}{Vadivel}
\title{GPU Software Engineer | High-Performance Computing | Scientific ML}
\email{srikrishnaa.careers@gmail.com}
\phone[mobile]{516-853-5663}
\social[linkedin]{sri-krishnaa-ece-phd}
\extrainfo{\href{https://krishnaa423.github.io/personal-website/}{Website} \textbar{} \href{https://github.com/krishnaa423}{GitHub: krishnaa423} \textbar{} \href{https://leetcode.com/u/krishnaa42342}{LeetCode: krishnaa42342} \textbar{} \href{https://hub.docker.com/u/krishnaa42342}{Docker Hub: krishnaa42342}}

\begin{document}
\makecvtitle

\section{Profile}
\cvitem{}{GPU-focused computational scientist and PhD researcher for performance-oriented systems and HPC roles. Develops Python and C++ code with CUDA, ROCm, MPI, and OpenMP, with experience in collective communication, multi-GPU parallelization, distributed debugging, and scaling analysis on national-lab HPC platforms.}

\section{Technical Skills}
\cvitem{Languages}{Python, C, C++, CUDA, JavaScript, Bash, PowerShell, SQL}
\cvitem{GPU / HPC}{CUDA, ROCm, MPI, OpenMP, PETSc, SLEPc, mpi4py, petsc4py, slepc4py, collective communication}
\cvitem{Platforms}{AMD MI250X, NVIDIA A100, Frontier, Perlmutter, Linux command line, macOS, Windows, CMake}
\cvitem{Debugging / Perf}{TotalView, Linaro DDT, CUDA-aware distributed debugging, strong scaling, weak scaling, hardware-aware optimization}
\cvitem{Machine Learning}{PyTorch, PyTorch Geometric, automatic differentiation, equivariant graph neural networks, scikit-learn}

\section{Research Experience}
\cventry{2021--present}{Research Assistant / PhD Researcher}{Yale University, Electrical Engineering}{}{}{\begin{itemize}%
\item Developed distributed scientific software in Python and C++ with MPI, OpenMP, PETSc, and SLEPc for large linear-algebra and simulation workloads as part of an INCITE award using exascale resources on Frontier at OLCF, Perlmutter at LBNL, and related national-lab HPC systems.
\item Parallelized matrix and vector workloads across multiple GPUs, using collective communication and distributed linear-algebra workflows to scale simulations across clusters.
\item Ported and optimized computational workloads for AMD MI250X and NVIDIA A100 GPU systems using CUDA and ROCm, translating many-body and excited-state physics methods into scalable GPU implementations.
\item Developed many-body and field-theoretic formulations of coupled electron, exciton, and phonon systems, including coherent-state and path-integral approaches for time-dependent and self-localized states.
\item Used PyTorch, PyTorch Geometric, and equivariant graph neural networks to accelerate self-trapped-exciton calculations and related materials simulations, achieving roughly 100x speedups in parts of the workflow by using training and automatic differentiation to obtain additional quantities without separate expensive calculations.
\item Ran strong- and weak-scaling studies and used TotalView, CUDA-aware MPI workflows, Linaro DDT, Bash, and Linux command-line tools to debug distributed applications across multi-process and multi-GPU environments.
\item Packaged scientific applications in Docker containers for deployment across different computing environments.
\item Co-authored a 2025 \textit{Journal of the American Chemical Society} paper and presented APS talks in 2024, 2025, and 2026 on self-trapped exciton formation, exciton-polarons, and related excited-state materials physics.
\end{itemize}}

\section{Professional Experience}
\cventry{2019--2021}{Software and Embedded Systems Engineer}{Probe Technologies Inc. / Weatherford Wireline Services}{}{}{\begin{itemize}%
\item Built CUDA-based data-processing pipelines and GPU-enabled engineering tooling for large acoustic waveform datasets generated by downhole logging tools.
\item Built backend tooling for database design, real-time remote tool testing, and calibration of temperature and pressure sensors.
\item Applied finite-element analysis with dolfinx for frequency analysis of steel pipes in well environments to support in-house physics and engineering studies.
\end{itemize}}
\cventry{2018--2019}{Photonics Technical Writer}{FindLight}{}{}{\begin{itemize}%
\item Produced technical articles and product-facing photonics content covering lasers, optics, and optoelectronic components for a specialized engineering readership.
\end{itemize}}

\section{Selected Projects}
\cventry{}{Multi-GPU Exciton-Phonon Simulation}{}{}{}{\begin{itemize}%
\item Developed distributed first-principles calculations for excited-state and exciton-phonon problems with CUDA, ROCm, MPI, and OpenMP, including multi-GPU matrix and vector parallelization on leadership-class GPU clusters.
\end{itemize}}
\cventry{}{Scientific ML on GPU Systems}{}{}{}{\begin{itemize}%
\item Developed PyTorch Geometric equivariant models for molecular-property prediction and for accelerating scientific calculations, with roughly 100x workflow speedups from combining training with automatic differentiation.
\end{itemize}}
\cventry{}{CUDA Containers for HPC Software}{}{}{}{\begin{itemize}%
\item Packaged scientific software in CUDA-ready Docker environments for portable deployment across local and HPC systems.
\end{itemize}}
\cventry{}{Matrix-Multiply AI Accelerator}{}{}{}{\begin{itemize}%
\item Built a Verilog-based matrix-multiply accelerator project for AI-oriented datapath design and hardware-software performance exploration.
\end{itemize}}
\section{Education}
\cventry{2021--present}{PhD, Electrical Engineering}{Yale University}{}{}{Ongoing doctoral research in computational materials, many-body theory, and accelerator-enabled scientific computing.}
\cventry{2023}{MPhil, Electrical Engineering}{Yale University}{}{}{}
\cventry{2023}{MS, Electrical Engineering}{Yale University}{}{}{}
\cventry{2017--2018}{MEng, Electrical and Computer Engineering}{Cornell University}{}{}{Included work in stochastic computing, device-oriented EE topics, and advanced design coursework.}
\cventry{2013--2017}{BS, Electrical and Computer Engineering}{Cornell University}{}{}{Included embedded systems, robotics, and core mathematics, physics, and computing foundations.}

\section{Selected Publications}
\cvitem{2025}{de Paula, A. M., Li, S., Hou, B., Vadivel, S., Teles-Ferreira, D. C., Iudica, A., Kabacinski, P., Hosseini, H., McArthur, J., Cerullo, G., et al. (2025). Time-domain observation of ultrafast self-trapped exciton formation in lead-free double halide perovskites. \textit{Journal of the American Chemical Society, 147}(32), 28923--28931.}

\section{Selected Conference Presentations}
\cvitem{2026}{Vadivel, S., Grande, R. D., Strubbe, D. A., \& Qiu, D. Y. (2026). First-principles study of exciton-polarons and self-trapped excitons. \textit{APS Global Physics Summit 2026}.}
\cvitem{2025}{Vadivel, S., Qiu, D. Y., Grande, R. D., \& Strubbe, D. A. (2025). First-principles study of self-trapped exciton formation in double perovskites. \textit{APS Global Physics Summit 2025}.}
\cvitem{2024}{Vadivel, S., \& Qiu, D. (2024). First-principles study of self-trapped exciton formation in double perovskites using excited state forces. \textit{APS March Meeting Abstracts 2024}, F01--007.}

\section{Selected Coursework}
\cvitem{Mathematics}{Differential Geometry; Abstract Algebra; Honors Mathematical Analysis II}
\cvitem{Computer Science}{Theoretical Computer Science; Probabilistic Machine Learning}
\cvitem{Physics}{Graduate Quantum Mechanics I--II; Solid State Physics I--II; Many-Body Quantum Physics}
\cvitem{Electrical Engineering}{Semiconductor Physics; Thin Film Science; Lasers and Optoelectronics; Integrated Photonics; VLSI Design; Analog VLSI Design; Mathematics of Signals and Systems; Advanced Embedded Systems}

\end{document}
